\chapter{Experiments}\label{ch:experiments}

\section{Experimental Protocol}\label{sec:protocol}
% Describe the overall objective of the experiments.
% The use of Leave-One-Out Cross-Validation (LOO-CV) over the healthy cohort.
% Why LOO-CV is appropriate given the limited number of independent healthy
% donors (patients Healthy_2 through Healthy_7).

\section{Leave-One-Out Cross-Validation Scheme}\label{sec:loo-cv}
% Detail the fold construction:
% - In each fold, hold out one healthy patient as the "test normal".
% - Train the OC-SVM on the remaining 5 healthy patients.
% - Evaluate the model on the held-out healthy patient and one tumour patient.
% - Repeat for each tumour patient (Colo_11, Colo_12, Colo_13) across all folds.
% Ensures that the test healthy patient is strictly unseen during training.

\section{Hyperparameter Selection}\label{sec:hyperparameters}
% Explain the fixed choices or search grid:
% - Kernel parameters: maximum k-mer size ($k_{max}$), allowed mismatches ($m$),
%   mismatch decay factor ($\lambda$).
% - MKL noise threshold.
% - OC-SVM hyperparameter: $\nu$ (expected anomaly rate).
% Mention the computational constraints that informed these choices.

\section{Implementation Details}\label{sec:experiment-details}
% Hardware used (e.g. CPU cores, memory limits).
% Sample sizes per patient (e.g. 18,000 for training, 1,500 for testing) to
% maintain manageable Gram matrix dimensions.
% Seeding and reproducibility.
